\documentclass[a4paper]{article}
\usepackage{amsfonts}
\usepackage{amsmath}
\usepackage{amssymb}
\usepackage{graphicx}     

\newcommand{\cis}{\operatorname{cis}}

\begin{document}
  
\noindent \large Auteur: Hendrik De Bruyn \\
\noindent \large Studierichting: Informatica\\
\noindent \large Oefeningenreeks 1C nr. 54\\

\medskip

\normalsize

\textbf{Opgave:} \\

Los op in $\mathbb{C}$: $z^6 + z^3 + 1 = 0$ \\

\textbf{Oplossing:} \\

Substitueer $v = z^3$. \\

Hieruit volgt dat $v^2 + v + 1 = 0$ \\

Discriminant: $D = 1 - 4 = -3$ \\

$-3 = (x + yi)^2 = x^2 - y^2 + 2xyi$ \\

$\Rightarrow\left\{\begin{array}{l}
    x^2 - y^2 = -3 \\
    2xy = 0
\end{array}\right.$ \\

$\Rightarrow -y^2 = -3 \Rightarrow y = \pm \sqrt{3}$ \\

$\Rightarrow x + yi = \sqrt{3}i$ en $x + yi = -\sqrt{3}i$ \\

$v_{1,2} = \dfrac{-1 \pm \sqrt{3}i}{2} = -\dfrac{1}{2} \pm \dfrac{\sqrt{3}}{2}i$ \\

Voor $v_1$ berekenen we de volgende drie oplossingen: \\

$-\dfrac{1}{2} + \dfrac{\sqrt{3}}{2}i = \cis \dfrac{2\pi}{3} = v_1 = z^3$ \\

$z = \cis \dfrac{\dfrac{2\pi}{3} + k\cdot2\pi}{3} = \cis \dfrac{2\pi + k\cdot6\pi}{9}$ (met $k \in \{0,1,2\})$ \\

${\bf z_0} = \cis \dfrac{2\pi}{9}$ \\

${\bf z_1} = \cis \dfrac{8\pi}{9}$ \\

${\bf z_2} = \cis \dfrac{14\pi}{9} = \cis \dfrac{-4\pi}{9}$ \\

Voor $v_2$ berekenen we de volgende drie oplossingen: \\

$-\dfrac{1}{2} - \dfrac{\sqrt{3}}{2}i = \cis \dfrac{4\pi}{3} = v_1 = z^3$ \\

$z = \cis \dfrac{\dfrac{4\pi}{3} + k\cdot2\pi}{3} = \cis \dfrac{4\pi + k\cdot6\pi}{9}$ (met $k \in \{0,1,2\})$ \\

${\bf z_3} = \cis \dfrac{4\pi}{9}$ \\

${\bf z_4} = \cis \dfrac{10\pi}{9} = \cis \dfrac{-8\pi}{9}$ \\

${\bf z_5} = \cis \dfrac{16\pi}{9} = \cis \dfrac{-2\pi}{9}$ \\

\begin{thebibliography}{999}
%\bibitem{a} Referentie
\end{thebibliography}
\end{document} 
